\Section{Подпоследовательности, частичные пределы. Теорема Больцано-Вейерштрасса.
Критерий Коши существования конечного предела последовательности.}

\Subsection{Подпоследовательности, частичные пределы}

\begin{definition}
Пусть \(\{x_{n}\}\)~--- произвольная последовательность,
а \(\{k_{n}\}\)~--- \textit{возрастающая} последовательность \textit{натуральных} чисел.
Тогда последовательность \(\{x_{k_{n}}\}\) называют \textbf{подпоследовательностью}
последовательности \(\{x_{n}\}\).
\end{definition}

\begin{note}
Заметим очевидную, но важную деталь~--- \(k_{n}\geq n\).
Тогда для фиксированного номера \(N\) и некоторого номера \(n\geq N\)
получаем 2 свойства:
\[
k_{n} \geq n,
\]
\[
k_{n} \geq N.
\]
\end{note}

Разберем свойства подпоследовательностей.

\begin{theorem} \label{theorem_subseries_1}
Если последовательность \(\{x_{n}\}\) сходится к числу \(\alpha\),
то и любая ее подпоследовательность также сходится к числу \(\alpha\).

\proofname~По определению предела:
\[
\forall \, \epsilon > 0 ~~ \exists \, N \in \N:
~ \forall n \in \N, ~ n \geq N \longmapsto |x_{n} - \alpha| < \epsilon.
\]
Пусть \(\{x_{k_{n}}\}\)~--- произвольная подпоследовательность исходной последовательности.
Тогда, согласно замечанию, \(k_{n} \geq N\), и выполняется \(|x_{k_{n}} - \alpha| < \epsilon\).
Следовательно, \(\{x_{k_{n}}\}\) сходится к \(\alpha\), что и требовалось доказать.~\qedsymbol

\begin{anote}
На всякий случай уточню, что отсчет в подпоследовательности мы ведем
не от номера \(N\), а от номера \(k_{N}\).
\end{anote}

\end{theorem}

\begin{theorem}
Если все подпоследовательности данной последовательности \(\{x_{n}\}\) сходятся,
то пределы этих последовательностей равны одному и тому же числу.

\vspace{2mm}

\proofname~Поскольку сама \(\{x_{n}\}\) является собственной подпоследовательностью,
то по условию \(\{x_{n}\}\) сходится. Тогда по Теореме \ref{theorem_subseries_1} пределы всех
подпоследовательностей равны пределу последовательности, что и требовалось доказать.~\qedsymbol
\end{theorem}

\begin{theorem} \label{theorem_BBP_series}
Каждая подпоследовательность ББП является ББП.

\proofname~По определению ББП для любого \(\gamma > 0\) существует
такой номер \(N\), что для всех номеров \(n \geq N\) выполняется \(|x_{n}| > \gamma\).
Поскольку \(k_{n} \geq N,\) то также выполняется \(|x_{k_{n}}| > \gamma\).
Следовательно, \(\{x_{k_{n}}\}\)~--- также ББП.~\qedsymbol
\end{theorem}

\begin{theorem}
Из любой сходящейся последовательности можно выделить
монотонную сходящуюся подпоследовательность.

\proofname~Пусть \(\lim\limits_{n \to \infty}{x_{n}} = \alpha\).
В произвольной \(\epsilon\)-окрестности выполняется \textbf{хотя бы}
одно из следующего:\\
\indent 1) содержится бесконечно много членов \(\{x_{n}\}\), равных \(\alpha\);\\
\indent 2) содержится бесконечно много членов \(\{x_{n}\}\), меньших \(\alpha\);\\
\indent 3) содержится бесконечно много членов \(\{x_{n}\}\), больших \(\alpha\).\\

\indent Если верно первое, то для доказательства достаточно составить подпоследовательность
из всех членов, равных \(\alpha\), расставив их по порядку. \\
\indent Рассмотрим случай 2. Он означает, что на интервале \((\alpha - \epsilon, ~ \alpha)\)
можно зафиксировать некоторый член \(x_{k_{1}}\). Теперь мы имеем два интервала:
\((\alpha - \epsilon, ~ x_{k_{1}}]\) и \((x_{k_{1}}, ~ \alpha)\).

Заметим, что левый интервал не может иметь бесконечно много членов~--- иначе из них можно было
бы составить подпоследовательность, не сходящуюся к \(\alpha\), что противоречило
бы Теореме \ref{theorem_subseries_1}. Следовательно, бесконечно много членов теперь содержит именно
правый интервал \((x_{k_{1}}, ~ \alpha)\). Перейдем к нему, и выберем на нем новый член
\(x_{k_{2}}\), причем \(k_{2} > k_{1}\) (мы можем это сделать т.к. членов бесконечно много).

Повторяя операцию, мы составим монотонную (т.к. \(x_{k_{n+1}} > x_{k_{n}}\))
подпоследовательность \(\{x_{k_{n}}\}\), сходящуюся к \(\alpha\)~---
что и требовалось доказать.~\qedsymbol
\end{theorem}

Теперь поговорим о частичных пределах.

\begin{definition}
Предел подпоследовательности \(\{x_{k_{n}}\}\) последовательности \(\{x_{n}\}\)
называется \textit{частичным пределом} последовательности \(\{x_{n}\}\).
\end{definition}

\begin{definition}
Число \(\alpha\) называется \textit{предельной точкой} последовательности \(\{x_{n}\}\),
если в любой ее \(\epsilon\)-окрестности содержится бесконечно
много членов последовательности \(\{x_{n}\}\).
\end{definition}

\begin{theorem}
Для того чтобы число \(\alpha\) являлось частичным пределом последовательности \(\{x_{n}\}\),
необходимо и достаточно, чтобы \(\alpha\) являлось предельной точкой этой последовательности.

\vspace{2mm}

\proofname~Докажем необходимость. \\
\indent Пусть \(\alpha\)~--- частичный предел \(\{x_{n}\}\), т.е. есть такая подпоследовательность
\(\{x_{k_{n}}\}\), что \(\lim\limits_{n \to \infty}{x_{k_{n}}} = \alpha.\)

По определению предела для любого положительного \(\epsilon\) существует такой номер \(N\),
что все члены \(x_{k_{n}}\) с номером, большим \(k_{N}\), будут находиться
в \(\epsilon\)-окрестности точки \(\alpha\). Причем заметим, что таких членов
будет бесконечно много (т.к. для \textbf{всех} \(n \geq N\)). Осталось добавить очевидное:
член последовательности \(\{x_{k_{n}}\}\) также является
членом последовательности \(\{x_{n}\}\). Тогда по определению
\(\alpha\) - предельная точка последовательности \(\{x_{n}\}\).\\
\\
\indent Теперь докажем достаточность. \\
\indent Пусть \(\alpha\)~--- предельная точка последовательности \(\{x_{n}\}\).
Тогда при \(\epsilon=1\) выберем номер \(k_{1}\) такой, что \(x_{k_{1}}\) принадлежит
интервалу \((\alpha - 1, ~ \alpha + 1)\). Затем при \(\epsilon = \frac{1}{2}\) выберем
такой номер \(k_{2}\), что \(k_{2} > k_{1}\) и член \(x_{k_{2}}\) принадлежит
интервалу \(\left( \alpha - \frac{1}{2}, ~ \alpha + \frac{1}{2} \right)\).

Будем повторять операцию до бесконечности, на \(n\)-ом шаге имея \(\epsilon = \frac{1}{n}\).
Заметим, что такая операция действительно возможна на каждом шаге, поскольку
у нас постоянно для выбора есть \textit{бесконечно много} членов последовательности.
Тогда мы построили подпоследовательность \(\{x_{k_{n}}\}\), сходящуюся к \(\alpha\)~---
то есть \(\alpha\) является частичным пределом для исходной последовательности \(\{x_{n}\}\),
что и требовалось доказать.~\qedsymbol
\end{theorem}

\begin{theorem}
Если последовательность \(\{x_{n}\}\) не ограничена сверху, то она содержит
бесконечно большую подпоследовательность \(\{x_{k_{n}}\}\) такую, что
\(\lim\limits_{n \to \infty}{x_{k_{n}}} = +\infty\).
Если последовательность \(\{x_{n}\}\) не ограничена снизу, то она содержит
бесконечно большую подпоследовательность \(\{x_{k_{n}}\}\) такую, что
\(\lim\limits_{n \to \infty}{x_{k_{n}}}= -\infty\).

\vspace{3mm}

\proofname~Докажем случай неограниченной сверху последовательности.\\
\indent Заметим важный факт: если мы отбросим конечное число членов неограниченной
сверху последовательности, то она останется неограниченной сверху.

Тогда проделаем следующую операцию: для числа \(\gamma = 1\) найдется такой
номер \(k_{1}\), что \(x_{k_{1}} > 1.\) \\

\indent Отбросим первые \(k_{1}\) членов последовательности. Для новой
(обратите внимание на границы!) последовательности \(\{x_{n}\}^{\infty}_{k_{1} + 1}\)
для \(\gamma=2\) найдется новый номер \(k_{2}\) такой, что \(x_{k_{2}} > 2\). \\

Повторяя операцию, мы определяем подпоследовательность \(\{x_{k_{n}}\}\),
в которой для всех номеров \(n\) выполняется \(x_{k_{n}} > n\).
По аксиоме Архимеда для любого \(\gamma > 0\) найдется натуральное число \(N > \gamma\) ---
следовательно, выполняется \(x_{k_{N}} > N > \gamma\). Для больших номеров (\(n \geq N\)) имеем:
\[
x_{k_{n}} > n \geq N > \gamma
\]

Следовательно, по определению \(\{x_{k_{n}}\}\) - ББП.~\qedsymbol
\end{theorem}

\Subsection{Теорема Больцано-Вейерштрасса}

\begin{theorem}[Теорема Больцано-Вейерштрасса]
Из любой ограниченной последовательности можно выделить сходящуюся подпоследовательность.

\vspace{3mm}

\proofname~Последовательность \(\{x_{n}\}\) ограничена~--- следовательно,
существуют такие числа \(\alpha\) и \(\beta\), что для любого номера \(n \in \N\)
выполняется \(\alpha \leq x_{n} \leq \beta\). \\

\indent Теперь будем делать следующее: разобьем отрезок \([\alpha, \, \beta]\) пополам.
Весь отрезок содержит бесконечно много членов последовательност~--- следовательно,
хотя бы одна из половинок обладает тем же свойством.
Тогда выберем ту половинку (или любую из двух подходящих), в которой осталось
бесконечно много членов, и обозначим ее \([\alpha_{1}, \, \beta_{1}]\).

Аналогичную операцию сделаем с отрезком \([\alpha_{1}, \, \beta_{1}]\)~---
получим отрезок \([\alpha_{2}, \, \beta_{2}]\) и так далее. \\
\indent Таким образом, мы определили систему отрезков \(\{[\alpha_{n}, \, \beta_{n}]\}\).

Заметим, что, во-первых, каждый следующий отрезок вложен в предыдущий,
а во-вторых, поскольку длина каждый раз уменьшается вдвое, имеем:
\[
\lim_{ n \to \infty } (\beta_{n} - \alpha_{n}) =
\lim_{ n \to \infty } \frac{\beta - \alpha}{2^n} = 0
\]

Тогда по определению система отрезков является стягивающейся.\\
\indent Отсюда по теореме Кантора существует единственная точка \(\gamma_{0}\), которая
принадлежит всем отрезкам системы. Осталось построить подпоследовательность,
сходящуюся к \(\gamma_{0}\).\\

\indent На отрезке \([\alpha_{1}, \, \beta_{1}]\) выберем первый по номеру член последовательности,
обозначив его номер как \(k_{1}\). Затем на отрезке \([\alpha_{2}, \, \beta_{2}]\)
выберем первый член последовательности с номером \(k_{2} > k_{1}\).
Заметим, что такая операция действительно возможна всегда,
потому что для выбора мы располагаем бесконечным количеством членов
исходной последовательности.\\

Продолжая операцию, получим подпоследовательность \(\{x_{k_{n}}\}\)
последовательности \(\{x_{n}\}\), сходящуюся к \(\gamma_{0}\).
Так мы получаем то, что и требовалось доказать.~\qedsymbol
\end{theorem}

\begin{theorem}[Теорема о единственном частичном пределе]
Если последовательность ограничена и имеет единственный частичный предел,
то она сходится к этому частичному пределу.

\vspace{3mm}

\proofname~Пусть для любого номера \(n\) выполняется \(\alpha \leq x_{n} \leq \beta\)
и \(\lim\limits_{n \to \infty}{x_{k_{n}}} = \gamma.\) Тогда по Теореме \ref{theorem_lim_compare}
(переход к пределам) имеем: \(\alpha \leq \gamma \leq \beta.\)

Предположим противное: пусть \(\gamma\) не является пределом исходной последовательности.
Тогда по отрицанию определения предела существует такое \(\epsilon_{0} > 0\),
что вне интервала \((\gamma - \epsilon_{0}, \, \gamma + \epsilon_{0})\) содержится
бесконечно много членов последовательности \(\{x_{n}\}\).
Тогда хотя бы с одной из сторон от интервала находится бесконечно много членов \(\{x_{n}\}\)
- допустим, что слева, то есть на отрезке \([\alpha, \, \gamma - \epsilon_{0}].\)
По теореме Больцано-Вейерштрасса на этом отрезке существует частичный предел,
отличный от \(\gamma\) - получаем противоречие с условием единственности.
Следовательно, предположение неверно, и теорема доказана.~\qedsymbol
\end{theorem}

\Subsection{Критерий Коши существования предела последовательности}

\begin{definition}[Вариант 1]
Числовая последовательность называется фундаментальной, если для любого
действительного числа \(\epsilon > 0\) найдется такой номер \(N\), что
для всех номеров \(n \geq N\) и \(m \geq N\) выполняется \(|x_{n} - x_{m}| < \epsilon\).
\[
\forall \, \epsilon > 0 ~~ \exists \, N : ~ \forall \, n \in \N, ~ n \geq N \hspace{1em}
\forall \, m \in \N, ~ m \geq N ~ \longmapsto ~ |x_{n} - x_{m}| < \epsilon
\]
\end{definition}

\begin{definition}[Вариант 2]
 Числовая последовательность называется фундаментальной, если для любого
действительного числа \(\epsilon > 0\) найдется такой номер \(N\), что
для всех номеров \(n \geq N\) и натурального \(p\) выполняется \(|x_{n+p} - x_{n}| < \epsilon\).
\[
\forall \, \epsilon > 0 ~~ \exists \, N : ~ \forall \, n \in \N, ~ n \geq N \hspace{1em}
\forall \, p \in \N ~ \longmapsto ~ |x_{n+p} - x_{n}| < \epsilon
\]
\end{definition}

\begin{lemma} \label{lemma_Koshi}
Если последовательность \(\{x_{n}\}\) фундаментальна, то она ограничена.

\vspace{3mm}

\proofname~Для фиксированного \(\epsilon = 1\) найдется номер \(N\) такой, что
для всех натуральных \(p\) выполняется неравенство \(|x_{N+p} - x_{N}| < 1\).
Следовательно, все члены последовательности \(\{x_{n}\}\) с номерами \(n \geq N\)
содержатся в интервале \((x_{N} - 1, \, x_{N} + 1)\). Заметим, что членов с номерами
меньшими, чем \(N\), \textit{конечное} количество. Тогда зададим следующее число:
\[
\gamma = \max(|x_{1}|, |x_{2}|, |x_{3}|, \dots, |x_{N-1}|, |x_{N}| + 1)
\]

Заметим, что теперь для всех номеров \(n\) выполняется \(|x_{n}| < \gamma\).
Следовательно, \(\{x_{n}\}\)~--- ограничена, что и требовалось доказать.~\qedsymbol
\end{lemma}

\begin{theorem}[Критерий Коши]
Для сходимости последовательности \(\{x_{n}\}\) необходимо и достаточно,
чтобы она была фундаментальной.

\vspace{3mm}

\proofname~Докажем необходимость.
Пусть \(\lim\limits_{n \to \infty} {x_{n}} = \alpha.\) По определению предела
для любого положительного \(\epsilon\) найдется такой номер \(N\), начиная с которого
 для всех номеров \(n \geq N\) будет выполняться \(|x_{n}-\alpha|<\frac{\epsilon}{2}.\)
Заметим, что для любых \(p \in \N\) также выполняется \(|x_{n+p} - \alpha| < \frac{\epsilon}{2}\).
Тогда для всех \(n \geq N\) и для любого \(p \in \N\) имеем:
\[
\begin{aligned}
|x_{n+p} - x_{n}| = |(x_{n+p} - \alpha) - (x_{n} - \alpha)| \leq \\
\leq |x_{n+p} - & \alpha| + |x_{n} - \alpha| < \frac{\epsilon}{2} + \frac{\epsilon}{2} = \epsilon
\end{aligned}
\]
То есть по определению \(x_{n}\)~--- фундаментальная. \\
\\

Теперь докажем достаточность.\\
\indent По Лемме \ref{lemma_Koshi} фундаментальная последовательность \(\{x_{n}\}\) ограничена.
Тогда по теореме Больцано-Вейерштрасса из исходной последовательности
можно выделить сходящуюся подпоследовательность \(\{x_{k_{n}}\}\). \\

\indent Пусть \(\lim\limits_{n \to \infty} {x_{k_{n}}} = \alpha.\)
Для произвольного \(\epsilon > 0\) по фундаментальности \(\{x_{n}\}\)
и сходимости \(\{x_{k_{n}}\}\) имеем:
\[
\exists \, N_{1} \in \N : ~ \forall n \in \N, ~ n \geq N_{1} \hspace{1em}
\forall m \in \N, ~ m \geq N_{1} ~ \longmapsto ~ |x_{n} - x_{m}| < \frac{\epsilon}{2},
\]
\[
\exists \, N_{2} \in \N : ~ \forall n \in \N, ~ n \geq N_{2}
~ \longmapsto ~ |x_{k_{n}} - \alpha| < \frac{\epsilon}{2}.
\]

Обозначим \(N = \max(N_{1}, N_{2}).\) Поскольку \(k_{N} \geq N,\) то имеем:
\[
\begin{aligned}
|x_{n} - \alpha| = |(x_{n} - x_{k_{N}}) + (x_{k_{N}} - \alpha)| \leq \\
\leq |x_{n} - x_{k_{N}}&| + |x_{k_{N}} - \alpha| <
\frac{\epsilon}{2} + \frac{\epsilon}{2} = \epsilon.
\end{aligned}
\]

Следовательно, по определению \(\alpha\)~--- предел \(\{x_{n}\}\),
то есть последовательность \(\{x_{n}\}\)~--- сходящаяся, что и требовалось доказать.~\qedsymbol
\end{theorem}
